\section{Метод решения}

Программа состоит из двух отдельных исполняемых файлов: родительского
процесса и дочернего процесса.

В начале работы родительская программа считывает с помощью стандартного
ввода имя файла, в который дочерняя программа должна записывать результаты.

Для передачи команд от родительского процесса дочернему создаётся
неименованный канал с помощью системного вызова \texttt{pipe()}.
Канал предоставляет два файловых дескриптора: один для чтения и один
для записи.

После создания канала родительский процесс вызывает \texttt{fork()}.
В результате создаётся дочерний процесс, наследующий открытые файловые
дескрипторы родителя.

В дочернем процессе закрывается дескриптор канала, предназначенный для
записи. Затем при помощи системного вызова

\begin{lstlisting}[language=C]
dup2(pipefd[0], STDIN_FILENO);
\end{lstlisting}

читающий конец канала связывается со стандартным входом дочернего
процесса.

Таким образом, после перенаправления стандартного входа дочерняя программа
может получать данные из канала при помощи обычных функций чтения из
\texttt{stdin}.

После настройки файловых дескрипторов дочерний процесс вызывает
\texttt{execl()}, заменяя выполняемую программу программой
\texttt{child}. Имя выходного файла передаётся дочерней программе в
качестве аргумента командной строки.

Схема взаимодействия процессов имеет следующий вид:

\begin{verbatim}
Пользователь
    |
    | имя файла
    v
Родительский процесс
    |
    | fork()
    +--------------------+
    |                    |
    |                Дочерний процесс
    |                    |
    |                dup2()
    |                    |
    |                exec(child)
    |                    |
    |                    v
    |                child
    |
    | команды
    | write()
    v
   pipe
    |
    v
 stdin дочернего процесса
    |
    v
 getline()
    |
    v
 strtof()
    |
    v
 вычисление суммы
    |
    v
 выходной файл
\end{verbatim}

Родительский процесс считывает команды пользователя с помощью
\texttt{getline()}. Полученная строка передаётся в канал функцией
\texttt{write\_all()}, которая повторяет вызов \texttt{write()} до тех
пор, пока не будут переданы все байты строки.

Использование дополнительной функции необходимо, поскольку один вызов
\texttt{write()} в общем случае может записать меньше байтов, чем было
запрошено. Также обрабатывается прерывание вызова сигналом
(\texttt{errno == EINTR}).

После окончания ввода родитель закрывает записывающий конец канала.
После этого дочерний процесс получает признак конца файла на стандартном
входе и завершает чтение.

Родительский процесс ожидает завершения дочернего процесса с помощью
\texttt{waitpid()}, что предотвращает появление процесса-зомби.

Второй канал для передачи данных от дочернего процесса родительскому в
данном варианте не требуется, поскольку результат должен записываться
непосредственно в файл.

\section{Описание программы}

Программа состоит из следующих исходных файлов:

\begin{itemize}
    \item \texttt{parent.c} --- создание дочернего процесса, создание
    канала, чтение пользовательского ввода и передача данных дочернему
    процессу;

    \item \texttt{child.c} --- обработка полученных строк, преобразование
    строковых представлений чисел в тип \texttt{float}, вычисление суммы
    и запись результатов в файл.
\end{itemize}

В родительской программе используются следующие основные системные
вызовы и функции POSIX:

\begin{itemize}
    \item \texttt{pipe()} --- создание неименованного канала;
    \item \texttt{fork()} --- создание дочернего процесса;
    \item \texttt{dup2()} --- перенаправление стандартного входа
    дочернего процесса на читающий конец канала;
    \item \texttt{execl()} --- замена программы дочернего процесса
    программой \texttt{child};
    \item \texttt{write()} --- передача данных через канал;
    \item \texttt{close()} --- закрытие неиспользуемых файловых
    дескрипторов;
    \item \texttt{waitpid()} --- ожидание завершения дочернего процесса.
\end{itemize}

Для чтения строк произвольной длины используется функция
\texttt{getline()}.

В дочерней программе функция \texttt{strtof()} используется для
последовательного преобразования чисел из текстового представления
в значения типа \texttt{float}. Указатель на первый символ после
распознанного числа позволяет последовательно обработать все числа
в строке независимо от их количества.

Выходной файл открывается дочерним процессом с помощью функции
\texttt{fopen()}. Для каждой полученной команды вычисленная сумма
записывается отдельной строкой с помощью \texttt{fprintf()}.

Ошибки системных вызовов проверяются по их возвращаемым значениям.
Для вывода текстового описания ошибки используется функция
\texttt{perror()}. Для функций преобразования чисел дополнительно
проверяется значение \texttt{errno}.